% diagrams/physical.tex — capa Physical: controles por anillos de defensa en profundidad,
% controles transversales y tiempos de recuperación (RTO/RPO).
\begin{tikzpicture}[x=1cm, y=1cm]
  % anillos de afuera hacia adentro: (gris de fondo, rótulo, controles)
  \foreach \i/\g/\capa/\ctl in {
      0/3/PROVEEDOR CLOUD/{certificaciones ISO 27001 y SOC 2 · región alterna},
      1/7/PERÍMETRO/{WAF · anti-DDoS · bots · mTLS al sistema de pago},
      2/11/RED/{VPC por dominio · subredes privadas · microsegmentación},
      3/15/HOST Y CONTENEDORES/{imágenes firmadas · hardening CIS · runtime},
      4/19/APLICACIÓN/{gateway OAuth2 · anti-replay · antifraude en línea · RASP}} {
    \pgfmathsetmacro{\T}{-0.47*\i}
    \pgfmathsetmacro{\B}{-5.15+0.36*\i}
    \pgfmathsetmacro{\L}{0.26*\i}
    \pgfmathsetmacro{\R}{9.3-0.26*\i}
    \draw[line width=0.7pt, fill=black!\g] (\L,\T) rectangle (\R,\B);
    \node[font=\tiny, anchor=north] at ({(\L+\R)/2},\T-0.05)
      {\textbf{\capa}\ \ \ctl};
  }
  % centro: los datos
  \node[fill=terminalaccent, text=white, font=\scriptsize\bfseries, align=center,
        minimum width=5.8cm, minimum height=1.05cm] (dat) at (4.65,-3.0)
    {DATOS\\[-1pt]{\mdseries\tiny cifrado con llaves por dominio · tokenización · llaves en HSM}};

  % el atacante entra desde afuera y cada anillo lo frena
  \draw[terminal arrow, draw=terminalaccent, line width=1pt] (-0.6,-2.9) -- (0.95,-2.9);
  \node[font=\tiny\bfseries, text=terminalaccent, anchor=east, align=right] at (-0.05,-2.55) {ataque};
  \node[font=\tiny, anchor=north, text=black!70] at (4.65,-5.2)
    {si un anillo cae, el siguiente contiene el ataque};

  % controles transversales: atraviesan todos los anillos
  \begin{scope}[every node/.style={terminal box, anchor=north west, text width=3.4cm,
                                   align=left, font=\tiny, inner sep=2.5pt}]
    \node (idn) at (9.75,0.0)
      {\textbf{IDENTIDAD}\\passkeys para clientes · llave FIDO2 para empleados · PAM con acceso temporal};
    \node (mon) at ([yshift=-7pt]idn.south west)
      {\textbf{MONITOREO}\\SIEM 24/7 · logs inmutables · alertas de fraude en vivo};
    \node (pro) at ([yshift=-7pt]mon.south west)
      {\textbf{PROCEDIMENTAL}\\segregación de funciones · doble aprobación · simulacros};
  \end{scope}
  \node[font=\tiny\bfseries, anchor=south west] at (9.75,0.02) {transversales};

  % tiempos de recuperación (RTO = cuánto puede estar caído; RPO = cuánto se puede perder)
  \node[terminal inv, anchor=north west, text width=3.4cm, align=left, font=\tiny\bfseries,
        inner sep=2.5pt] (th) at ([yshift=-9pt]pro.south west) {RTO / RPO};
  \node[draw, line width=0.6pt, anchor=north west, text width=3.4cm, align=left,
        font=\tiny, inner sep=2.5pt] at (th.south west)
    {\textcolor{terminalaccent}{\textbf{pagos: RPO 0}}\\
     RTO < 1 min (zona) · ≤ 2 h (región)\\[1pt]
     llaves y QR: RTO 1 h · RPO 5 min\\[1pt]
     backoffice: RTO y RPO 24 h};
\end{tikzpicture}
